\documentclass[12pt]{article}
\usepackage[a4paper,margin=1in]{geometry}
\usepackage[utf8]{inputenc}
\usepackage[T1]{fontenc}
\usepackage{amsmath,amssymb}
\usepackage{textcomp}
\usepackage{parskip}
\begin{document}
\begin{center}
{\Large \textbf{2008 STEP I}}
\end{center}

\subsection*{Page 1}
\begin{flushleft}
\bigskip\textbf{Section A: Pure Mathematics}\par
1 What does it mean to say that a number x is irrational?\par
Prove by contradiction statements A and B below, where p and q are real numbers.\par
A: If pq is irrational, then at least one of p and q is irrational.\par
B: If p + q is irrational, then at least one of p and q is irrational.\par
Disprove by means of a counterexample statement C below, where p and q are real numbers.\par
C: If p and q are irrational, then p + q is irrational.\par
2 2\par
If the numbers e, \pi, \pi , e and e\pi are irrational, prove that at most one of the numbers \pi + e,\par
2 2 2 2\par
\[\pi e, \pi e , \pi + e is rational.\]
- -\par
\bigskip\textbf{Question 2.} The variables t and x are related by t = x + \sqrtx2 + 2bx + c , where b and c are constants and\par
\[b < c. Show that\]
dx t x\par
\[= - ,\]
dt t + b\par
and hence integrate .\par
\[\sqrt\{x^2\} + 2bx + c\]
Verify by direct integration that your result holds also in the case b = c if x + b > 0 but that\par
\[your result does not hold in the case b = c if x + b < 0 .\]
\end{flushleft}
\subsection*{Page 2}
\begin{flushleft}
\bigskip\textbf{Question 3.} Prove that, if c > a and d > b, then\par
\[ab + cd > bc + ad . ( )\]
∗\par
\textbf{(i)} If x > y, use ( ) to show that x 2 + y 2 > 2xy .\par
∗\par
\[If, further, x > z and y > z, use ( ) to show that z 2 + xy > xz + yz and deduce that\]
\[x 2 + y 2 + z 2 > xy + yz + zx . ∗\]
\[Prove that the inequality x 2 + y 2 + z 2 > xy + yz + zx holds for all x, y and z.\]
\textbf{(ii)} Show similarly that the inequality\par
s t r\par
\[+ + > 3\]
t r s\par
holds for all positive r, s and t.\par
[Note: The final part of this question differs (though not substantially) from what appeared\par
in the actual examination since this was found to be unsatisfactory (though not incorrect) in\par
a way that had not been anticipated.]\par
\bigskip\textbf{Question 4.} A function f(x) is said to be convex in the interval a < x < b if f 00 (x) > 0 for all x in this\par
interval.\par
2 2\par
\textbf{(i)} Sketch on the same axes the graphs of y = cos x and y = sin x in the interval\par
\[0 6 x 6 2\pi.\]
\[The function f(x) is defined for 0 < x < 2\pi by\]
2 sin x\par
\[f(x) = e 3 .\]
Determine the intervals in which f(x) is convex.\par
\textbf{(ii)} The function g(x) is defined for 0 < x < \pi by\par
-k tan x\par
\[g(x) = e .\]
\[If k = sin 2α and 0 < α < \pi/4, show that g(x) is convex in the interval 0 < x < α, and\]
give one other interval in which g(x) is convex.\par
\end{flushleft}
\subsection*{Page 3}
\begin{flushleft}
\bigskip\textbf{Question 5.} The polynomial p(x) is given by\par
n-1\par
n r\par
\[p(x) = x + arx ,\]
\[r=0\]
P\par
where a0, a1, . . . , an-1 are fixed real numbers and n > 1. Let M be the greatest value of\par
\[p(x) for x 6 1. Then Chebyshev's theorem states that M > 2 1-n .\]
| |\par
\{ \{\par
(\{ i) \{Prove Chebyshev's theorem in the case n = 1 and verify that Chebyshev's theorem holds\par
in the following cases:\par
2 1\par
\[(a) p(x) = x ;\]
-\par
\[(b) p(x) = x x .\]
-\par
5 4 3 2\par
\textbf{(ii)} Use Chebyshev's theorem to show that the curve y = 64x + 25x 66x 24x + 3x + 1\par
- -\par
has at least one turning point in the interval 1 6 x 6 1.\par
-\par
\bigskip\textbf{Question 6.} The function f is defined by\par
x\par
e 1\par
\[f(x) = - , x > 0,\]
e 1\par
-\par
and the function g is the inverse function to f, so that g(f(x)) = x. Sketch f(x) and g(x) on\par
the same axes.\par
Verify, by evaluating each integral, that\par
1 k\par
2 1\par
\[f(x) dx + g(x) dx = ,\]
\[0 0 2(\sqrt\{e\} + 1)\]
Z Z\par
\[where k = , and explain this result by means of a diagram.\]
\[\sqrt\{e\} + 1\]
\bigskip\textbf{Question 7.} The point P has coordinates (x, y) with respect to the origin O. By writing x = r cos θ and\par
\textasciicircum{}\circ\par
\[y = r sin θ, or otherwise, show that, if the line OP is rotated by 60 clockwise about O,\]
the new y-coordinate of P is 1 (y \sqrt3 x). What is the new y-coordinate in the case of an\par
\textasciicircum{}\circ -\par
anti-clockwise rotation by 60 ?\par
An equilateral triangle OBC has vertices at O, (1, 0) and ( 1 , 1 \sqrt3), respectively. The point P\par
2 2\par
has coordinates (x, y). The perpendicular distance from P to the line through C and O is h1;\par
\[the perpendicular distance from P to the line through O and B is h^2; and the perpendicular\]
\[distance from P to the line through B and C is h^3.\]
\[Show that h^1 =\]
\[1 y \sqrt\{3\} x and find expressions for h^2 and h^3.\]
-\par
Show that h1 + h2 \{ \{ + h3 = 1 2 \{ \{ \sqrt3 if and only if P lies on or in the triangle OBC.\par
\end{flushleft}
\subsection*{Page 4}
\begin{flushleft}
\bigskip\textbf{Question 8.} (i) The gradient y of a curve at a point (x, y) satisfies\par
0 2 0\par
\[(y ) xy + y = 0 . ( )\]
- ∗\par
00 0\par
\[By differentiating ( ) with respect to x, show that either y = 0 or 2y = x .\]
∗\par
\[Hence show that the curve is either a straight line of the form y = mx + c, where\]
2 2\par
\[c = m , or the parabola 4y = x .\]
-\par
\textbf{(ii)} The gradient y of a curve at a point (x, y) satisfies\par
2 0 2 0 2\par
\[(x 1)(y ) 2xyy + y 1 = 0 .\]
- - -\par
Show that the curve is either a straight line, the form of which you should specify, or a\par
circle, the equation of which you should determine.\par
\end{flushleft}
\subsection*{Page 5}
\begin{flushleft}
\bigskip\textbf{Section B: Mechanics}\par
\bigskip\textbf{Question 9.} Two identical particles P and Q, each of mass m, are attached to the ends of a diameter of a\par
light thin circular hoop of radius a. The hoop rolls without slipping along a straight line on a\par
horizontal table with the plane of the hoop vertical. Initially, P is in contact with the table.\par
At time t, the hoop has rotated through an angle θ. Write down the position at time t of P ,\par
relative to its starting point, in cartesian coordinates, and determine its speed in terms of a,\par
θ and θ˙. Show that the total kinetic energy of the two particles is 2ma 2 θ˙2 .\par
Given that the only external forces on the system are gravity and the vertical reaction of the\par
table on the hoop, show that the hoop rolls with constant speed.\par
\bigskip\textbf{Question 10.} On the (flat) planet Zog, the acceleration due to gravity is g up to height h above the surface\par
and g at greater heights. A particle is projected from the surface at speed V and at an angle\par
2 2\par
\[α to the surface, where V sin α > 2gh . Sketch, on the same axes, the trajectories in the\]
0 0\par
\[cases g = g and g < g.\]
Show that the particle lands a distance d from the point of projection given by\par
0 0\par
V V V\par
\[d = - + V sin 2α ,\]
\[g g^0\]
\lambda \mu\par
\[where V 0 = V 2 2gh cosec2α .\]
-\par
p\par
\bigskip\textbf{Question 11.} A straight uniform rod has mass m. Its ends P1 and P2 are attached to small light rings that\par
are constrained to move on a rough rigid circular wire with centre O fixed in a vertical plane,\par
\[and the angle P1OP2 is a right angle. The rod rests with P^1 lower than P^2, and with both\]
ends lower than O. The coefficient of friction between each of the rings and the wire is µ.\par
Given that the rod is in limiting equilibrium (i.e. on the point of slipping at both ends), show\par
that\par
1 2µ µ\par
\[tan α = - - ,\]
1 + 2µ µ2\par
-\par
\textasciicircum{}\circ\par
\[where α is the angle between P1O and the vertical (0 < α < 45 ).\]
\[Let θ be the acute angle between the rod and the horizontal. Show that θ = 2λ, where λ is\]
\textasciicircum{}\circ\par
\[defined by tan λ = µ and 0 < λ < 22.5 .\]
\end{flushleft}
\subsection*{Page 6}
\begin{flushleft}
\bigskip\textbf{Section C: Probability and Statistics}\par
\bigskip\textbf{Question 12.} In this question, you may use without proof the results:\par
n n\par
1 2 1\par
\[r = n(n + 1) and r = n(n + 1)(2n + 1) .\]
2 6\par
\[r=1 r=1\]
X X\par
The independent random variables X1 and X2 each take values 1, 2, . . ., N , each value being\par
equally likely. The random variable X is defined by\par
\[X^1 if X^1 > X^2\]
\[X =\]
\[(X^2 if X^2 > X^1 .\]
2r 1\par
\textbf{(i)} Show that P(X = r) = - for r = 1, 2, . . ., N .\par
N 2\par
\textbf{(ii)} Find an expression for the expectation, µ, of X and show that µ = 67.165 in the case\par
\[N = 100.\]
\textbf{(iii)} The median, m, of X is defined to be the integer such that P(X > m) > 1 and\par
\[P(X 6 m) > 1 . Find an expression for m in terms of N and give an explicit value for\]
\[m in the case N = 100.\]
\textbf{(iv)} Show that when N is very large,\par
\[µ 2\sqrt\{2\}\]
.\par
m \approx 3\par
\bigskip\textbf{Question 13.} Three married couples sit down at a round table at which there are six chairs. All of the\par
possible seating arrangements of the six people are equally likely.\par
\textbf{(i)} Show that the probability that each husband sits next to his wife is .\par
\textbf{(ii)} Find the probability that exactly two husbands sit next to their wives.\par
\textbf{(iii)} Find the probability that no husband sits next to his wife.\par
\end{flushleft}
\end{document}
